\tcbset{colframe = blue!50!black, 
colback = white, 
toptitle=1mm, 
bottomtitle=1mm, 
fonttitle=\fontsize{10pt}{12pt}\selectfont,  breakable, enhanced jigsaw, finish broken={\tcbset{colframe=white, bottom=0mm, bottomrule=0mm}},before upper={\tcbset{colframe=white, toprule=0mm}},after={\tcbset{colframe=blue!40!black, toprule=1mm, bottomrule=1mm}},}

\section{Details of \ours Environment}

\subsection{Environment Infrastructure}
\label{app:env_infra}
As compared to core commonly used techniques like Docker~\footnote{\url{https://www.docker.com/}}, virtual machines can operate their own kernel and system, enabling compatibility with a wide variety of operating systems (such as Windows, macOS, Linux, etc.) across different CPU hardware types (x64, ARM, etc.), and supports training and evaluation in a multiprocess manner on both headless servers and personal computers.

\subsection{Observation Space}
\label{app:observation_space}
We implement three kinds of observation: 
\textbf{complete screenshot}, \textbf{accessibility tree} and \textbf{terminal output}.
We also implement a video recorder of the environment but don't put it into our modeling due to the agent's ability limitations.
\ours supports observation refactoring and extending if needed, such as getting data from certain opening applications that we want to focus on.

\subsubsection{Screenshot}
To align with the perception of a human user, we capture a screenshot of the entire computer screen. 
Including the mouse cursor also proves helpful in certain cases where mouse information is crucial. 
For screen resolution, we default to 1920$\times$1080, as it is the most commonly used screen resolution according to Internet Users Screen Resolution Realtime Statistics for 2023\footnote{\url{https://www.screenresolution.org/year-2023/}}. 
This resolution also offers a 16:9 aspect ratio. 
\ours also supports modifying the resolution of virtual machines to avoid potential memorization of absolute pixel values and to assist studies on topics like generalization across different resolutions.

% \begin{figure}[H]
%     \centering
%     \includegraphics[width=0.8\textwidth]{images/environment/desktop_screenshot.png}
%     \caption{Example of the screenshot with the mouse cursor captured on the desktop.}
%     \label{fig:desktop_screenshot_example}
% \end{figure}


\subsubsection{Accessibility Tree}
An accessibility tree (or a11y tree, same logic to kubernetes and k8s), refers to an intricate structure generated by the browser or OS accessibility APIs that renders a representative model of the web content, providing a means of interaction for assistive technologies. 
Each node within the accessibility tree hosts important information about a UI element. 
This could range from the nature of the object (a button, checkbox, or paragraph of text), its current state (checked or unchecked, for checkboxes), and even its spatial orientation on the screen.

Different operating systems employ varied accessibility APIs and tools to construct and manipulate the accessibility tree. 
These include Microsoft Active Accessibility (MSAA) and User Interface Automation (UIA) for Windows, NSAccessibility Protocol and macOS Accessibility Inspector for macOS, and 
Assistive Technology Service Provider Interface (ATSPI)~\footnote{\url{https://docs.gtk.org/atspi2/}} for GNOME dekstop used on Ubuntu.
We adopt \texttt{pyatspi} to get the accessibility tree on Ubuntu and \verb|pywinauto| on Windows.
We further convert it into XML format for message passing.
Partial pieces of the XML formatted accessibility tree are shown in Figure~\ref{fig:accessibility_tree_xml}.
\begin{figure}[ht]
    \centering
    \includegraphics[width=0.8\textwidth]{images/environment/accessibility_tree_xml.png}
    \caption{Raw XML formatted accessibility tree visualization.}
    \label{fig:accessibility_tree_xml}
\end{figure}


Tools such as \texttt{Accerciser} can help visualize the corresponding relationship of tree nodes and GUI components in the accessibility tree as shown in Figure~\ref{fig:accerciser_example}.

\begin{figure}[htbp]
    \centering
    \includegraphics[width=0.8\textwidth]{images/environment/accerciser.png}
    \caption{Use \texttt{Accerciser} to visualize components in accessibility tree. This example shows the corresponding relationship between the ``Print'' button (which is a printer icon on the bar of the LibreOffice app as shown on the left) and the ``Print'' tree node (which is the ``Print'' of the Accerciser Accessibility Explorer as shown on the right).}
    \label{fig:accerciser_example}
\end{figure}


\subsection{Action Space}
\label{app:action_space}
We implement two kinds of action space: 
\texttt{pyautogui} and \texttt{computer\_13}.
We mainly use \texttt{pyautogui} action space, since it saves tokens for describing action space definition in prompting, compared with self-designed actions.

\subsubsection{\texttt{PYAUTOGUI}}
\texttt{pyautogui} is an open-source, cross-platform Python module utilized for programmatically controlling the mouse and keyboard. 
It can control simple movements, clicks, as well as keyboard inputs, and can provide screen capture features or locate where the screen-specific images reside that can be useful for GUI automation. 
Compared with other Python packages such as xdotool~\footnote{\url{https://github.com/jordansissel/xdotool}}, mouse~\footnote{\url{https://github.com/boppreh/mouse}} and keyboard~\footnote{\url{https://github.com/boppreh/keyboard}}, \texttt{pyautogui} has better compatibility across different OSes as found in our early attempts, so that we adapt that as our core component of computer controlling as well as an official valid action space.

Herein, we will demonstrate some use cases of \texttt{pyautogui} and illustrate how it can be wielded as an action space.

\paragraph{Mouse Controlling Functions}
\small
\begin{Verbatim}
import pyautogui

# move mouse to XY coordinates over num_second seconds
pyautogui.moveTo(x, y, duration=num_seconds)

# drag mouse to XY
pyautogui.dragTo(x, y, duration=num_seconds)

pyautogui.click(x=moveToX, y=moveToY, clicks=num_of_clicks, button='left')
pyautogui.rightClick(x=moveToX, y=moveToY)
pyautogui.middleClick(x=moveToX, y=moveToY)
pyautogui.doubleClick(x=moveToX, y=moveToY)
pyautogui.tripleClick(x=moveToX, y=moveToY)
pyautogui.scroll(amount_to_scroll, x=moveToX, y=moveToY)
pyautogui.mouseDown(x=moveToX, y=moveToY, button='left')
pyautogui.mouseUp(x=moveToX, y=moveToY, button='left')
\end{Verbatim}

\paragraph{Keyboard Controlling Functions}

\small
\begin{Verbatim}
import pyautogui

# useful for entering text, newline is Enter
pyautogui.typewrite('Hello world!\n', interval=secs_between_keys)

pyautogui.typewrite(['a', 'b', 'c', 'left', 'backspace', 'enter', 'f1'], interval=secs_between_keys)
pyautogui.hotkey('ctrl', 'c')  # ctrl-c to copy
pyautogui.hotkey('ctrl', 'v')  # ctrl-v to paste
pyautogui.keyDown(key_name)
pyautogui.keyUp(key_name)
\end{Verbatim}

\paragraph{\texttt{pyautogui} as an Action Space}

Given the various controls it provides, \texttt{pyautogui} can readily be used as an action space in building automation software or testing interfaces with minor adjustments.
More formally, an action is within the action space when it meets the syntax of \texttt{pyautogui} or is one of three special actions \texttt{WAIT}, \texttt{FAIL}, and \texttt{DONE}.
This might include actions like clicking at a certain location, entering text or key sequences, or even resting for a span (Pause). Each action could be mapped to an identifying label or number, forming a discrete action space. 
For example:
\small
\begin{Verbatim}
import pyautogui

def perform_action(action):
  if action == 0:
    pyautogui.moveTo(100, 100)
  elif action == 1:
    pyautogui.write('Hello world!', interval=0.25)
  else:
    pyautogui.pause(1)
\end{Verbatim}
In this scheme, the "perform\_action" function constitutes the action space, where each unique action is associated with a unique integer (its action ID). 
The function interprets these action IDs and performs the corresponding action, forming a rudimentary discrete action space.

One interesting finding is that language models generate screenshot locate functions like:
\begin{Verbatim}
pyautogui.locateOnScreen('Apple.png')
\end{Verbatim}
When there is insufficient grounding evidence (such as when no screenshot is inputted, the accessibility tree lacks a specific element, or the multimodal model cannot comprehend the user interface), employing this function to retrieve the correct icon image could present an interesting method.

\subsubsection{\texttt{COMPUTER\_13}}
To facilitate potential reinforcement learning applications, we have created a variant of \texttt{pyautogui}, which we named \texttt{computer\_13}. 
In this variant, we wrap \texttt{pyautogui} into a finite action class with parameterized enumeration, such that it features 13 action types, excluding three special ones for task process control. 
Utilizing this structured approach allows more effective reinforcement learning by providing a distinct and finite set of actions to be learned and optimized.
As summarized in Table~\ref{table:action_space}, each action type has certain parameters, detailed in the collection, confirming the type, range, and whether each parameter is optional for that action.

\begin{table}[ht]
\centering
\caption{Action types and parameters defined in action space \texttt{computer\_13}, a variance we created for the potential reinforcement learning research based on our environment.}
\label{table:action_space}
\begin{tabular}{p{2.4cm}|p{2cm}|p{8.5cm}}
\textbf{Action Type} & \textbf{Parameters} & \textbf{Note} \\
\hline
\texttt{MOVE\_TO} & \textit{x, y} & \textit{Move the cursor to the specified position} \\
\texttt{CLICK} & \textit{button, \newline x, y, \newline num\_clicks} & \textit{Click the left button if the button not specified, otherwise click the specified button; click at the current position if x and y are not specified, otherwise click at the specified position} \\
\texttt{MOUSE\_DOWN} & \textit{button} & \textit{Press the left button if the button not specified, otherwise press the specified button} \\
\texttt{MOUSE\_UP} & \textit{button} & \textit{Release the left button if the button not specified, otherwise release the specified button} \\
\texttt{RIGHT\_CLICK} & \textit{x, y} & \textit{Right click at the current position if x and y are not specified, otherwise right click at the specified position} \\
\texttt{DOUBLE\_CLICK} & \textit{x, y} & \textit{Double click at the current position if x and y are not specified, otherwise double click at the specified position} \\
\texttt{DRAG\_TO} & \textit{x, y} & \textit{Drag the cursor to the specified position with the left button pressed} \\
\texttt{SCROLL} & \textit{dx, dy} & \textit{Scroll the mouse wheel up or down} \\
\texttt{TYPING} & \textit{text} & \textit{Type the specified text} \\
\texttt{PRESS} & \textit{key} & \textit{Press the specified key and release it} \\
\texttt{KEY\_DOWN} & \textit{key} & \textit{Press the specified key} \\
\texttt{KEY\_UP} & \textit{key} & \textit{Release the specified key} \\
\texttt{HOTKEY} & \textit{keys} & \textit{Press the specified key combination} \\
\hline
\texttt{WAIT} & - & \textit{Wait until the next action} \\
\texttt{FAIL} & - & \textit{Decide the task cannot be performed} \\
\texttt{DONE} & - & \textit{Decide the task is done} \\
\end{tabular}
\end{table}


\section{Details of \ours Benchmark}
\subsection{Operating System Selection}
\label{app:os_selection}
As of 2023, the most popular desktop operating systems are Windows (69.5\%), macOS (20.4\%), ChromeOS (3.2\%), and Linux (3.1\%) \footnote{\url{https://www.statista.com/statistics/218089/global-market-share-of-windows-7/}}. 
While Windows and macOS dominate the market share and boast the richest software ecosystems, their closed-source nature raises potential copyright concerns for direct usage. 
ChromeOS, being a web-based operating system, heavily depends on a Google account for its functionalities, rendering it less suitable for a public benchmark. 

In contrast, Linux desktops offer a wealth of open-source software for most daily tasks, supported by an active community for both basic and advanced use. 
Essential applications such as Office Suite, browsers, email clients, multimedia apps, and thousands of games and applications are either pre-installed or readily available through the software center of the distribution. 
Consequently, we select Ubuntu, the most representative Linux desktop OS, as the foundation for the main part of our benchmark intended for public use.
Additionally, we have developed components to facilitate agent interaction on a Windows virtual machine and have created a relatively small set of examples focusing on the Microsoft Office suite, including Excel, PowerPoint, and Word. 
This serves as a counterpart to the LibreOffice suite available on Ubuntu. 
These components can be utilized in-house or officially with the purchase of a license.
Regarding macOS, theoretically, it is illegal to install macOS on non-Apple devices, which leads us to refrain from developing our benchmark on this platform to avoid copyright issues.

\subsection{Software Selection}
\label{app:software_selection}
Due to the high cost of obtaining operation and evaluation script annotation data, we have chosen a representative set of software for the examples of Ubuntu part.
We adopt standards that consider: 
1) Availability - the software must be available on Ubuntu 22.04; 
2) Open-source - the software should be open-sourced with an appropriate license to prevent copyright issues; 
3) Popularity - the software should take a high download number and frequency of recommendations in blogs and tutorials; 
4) Strong user community and good support resources - it is preferable to have an active and robust user community as well as official documents, which can serve as ample resources for task collection and agent learning; 
5) Diversity of categories - the software should be diverse to adequately represent and cover a wide range of real-world cases. 

As a result, we have shortlisted the software into two categories: general usage and professional usage. 
For general usage, we have VLC for media playback, Chrome for web browsing, and Thunderbird for email management.
For professional usage, we have VS Code as a coding IDE, and LibreOffice (Calc, Writer, and Impress) for handling spreadsheets, documents, and presentations respectively, along with GIMP for image editing. 
This brings our total to eight different types of software.

\subsection{Task Example Sources}
\label{sub:task_exam_srcs}

We detail the task example sources in Table~\ref{tab:task_exam_srcs}.

\begin{longtable}{cll}
    \caption{Task Example Resources} \label{tab:task_exam_srcs} \\
\textbf{App} & \textbf{Resources} & \textbf{Link} \\
\hline
\endfirsthead

\multicolumn{2}{c}%
{{\bfseries \tablename\ \thetable{} -- continued from previous page}} \\
\hline
\textbf{App} & \textbf{Resources} \\
\hline
\endhead

\multicolumn{2}{r}{{\textit{Continued on next page}}} \\ 
\endfoot

\hline
\endlastfoot

\multirow{5}{*}{OS} &    Ubuntu Documentations & \url{https://help.ubuntu.com/} \\
&   Ask Ubuntu & \url{https://askubuntu.com/} \\
&   Super User & \url{https://superuser.com/} \\
&   Stack Overflow & \url{https://stackoverflow.com} \\
&   YouTube & \url{https://www.youtube.com/} \\
\hline
\multirow{10}{*}{Calc} &   LibreOffice Help & \url{https://help.libreoffice.org/} \\
&   Microsoft Tech Community & \url{https://techcommunity.microsoft.com/} \\
&   libreofficehelp.com & \url{https://www.libreofficehelp.com/} \\
&   Reddit r/LibreOfficeCal & \url{https://www.reddit.com/r/LibreOfficeCalc/} \\
&   Reddit r/Excel & \url{https://www.reddit.com/r/Excel/} \\
&   Super User & \url{https://superuser.com/} \\
&   Medium & \url{https://medium.com/} \\
&   Quora & \url{https://www.quora.com/} \\
&   YouTube & \url{https://www.youtube.com/} \\
&   Ryan and Debi \& Toren Personal Site & \url{https://www.ryananddebi.com/} \\
\hline
\multirow{9}{*}{Writer} &   LibreOffice Help & \url{https://help.libreoffice.org/} \\
&   LibreOffice Forum & \url{https://ask.libreoffice.org/} \\
&   libreofficehelp.com & \url{https://www.libreofficehelp.com/} \\
&   Super User & \url{https://superuser.com/} \\
&   Stack Overflow & \url{https://stackoverflow.com} \\
&   Ask Ubuntu & \url{https://askubuntu.com/} \\
&   Quora & \url{https://www.quora.com/} \\
&   YouTube & \url{https://www.youtube.com/} \\
&   SearchStar Personal Site & \url{https://seekstar.github.io/} \\
\hline
\multirow{8}{*}{Impress} & LibreOffice Help & \url{https://help.libreoffice.org/} \\
&   libreofficehelp.com & \url{https://www.libreofficehelp.com/} \\
&   Reddit r/LibreOffice & \url{https://www.reddit.com/r/LibreOffice/} \\
& Super User & \url{https://superuser.com/} \\
&   Stack Overflow & \url{https://stackoverflow.com} \\
& Technical Tips & \url{https://technical-tips.com/} \\
&   Just Click Here & \url{https://justclickhere.co.uk/} \\
&   TikTok & \url{https://www.tiktok.com/} \\
\hline
\multirow{9}{*}{VLC} &  VLC Documentation & \url{https://docs.videolan.me} \\
&   VLCHelp.com & \url{https://www.vlchelp.com/} \\
&   VideoLAN's Wiki & \url{https://wiki.videolan.org/} \\
&   Ubuntu Documentations & \url{https://help.ubuntu.com/} \\
&   Reddit r/Fedora & \url{https://www.reddit.com/r/Fedora/}\\
&   Super User & \url{https://superuser.com/} \\
&   Medium & \url{https://medium.com/} \\
&   YouTube & \url{https://www.youtube.com/} \\
&   Dedoimedo & \url{https://www.dedoimedo.com/index.html} \\
\hline
\multirow{8}{*}{Thunderbird} &  Thunderbird Support & \url{https://support.mozilla.org/en-US/products/thunderbird}\\
&   Reddit r/Thunderbird & \url{https://www.reddit.com/r/Thunderbird/} \\
&   Reddit r/Automation & \url{https://www.reddit.com/r/automation/} \\
&   Super User & \url{https://superuser.com/} \\
&   WikiHow & \url{https://www.wikihow.com/} \\
&   Quora & \url{https://www.quora.com/} \\
&   BitRecover & \url{https://www.bitrecover.com/} \\
&   AdSigner & \url{https://www.adsigner.com/} \\
\hline
Chrome &    Google Chrome Help & \url{https://support.google.com/chrome} \\
&   Reddit r/Chrome & \url{https://www.reddit.com/r/Chrome/} \\
&   Super User & \url{https://superuser.com/} \\
&   WikiHow & \url{https://www.wikihow.com/} \\
&   in5steps.com & \url{https://in5stepstutorials.com/} \\
&   How-To Geek & \url{https://www.howtogeek.com/} \\
&   Medium & \url{https://medium.com/} \\
&   Quora & \url{https://www.quora.com/} \\
&   YouTube & \url{https://www.youtube.com/} \\
&   Laptop Mag & \url{https://www.laptopmag.com} \\
\hline
\multirow{5}{*}{VS Code} &   Super User & \url{https://superuser.com/} \\
&   Stack Overflow & \url{https://stackoverflow.com} \\
& Quora & \url{https://www.quora.com/} \\
&   YouTube & \url{https://www.youtube.com/} \\
&   Campbell Muscle Lab GitHub & \url{https://campbell-muscle-lab.github.io/} \\
\hline
\multirow{6}{*}{GIMP} &   Reddit r/GIMP & \url{https://www.reddit.com/r/GIMP/} \\
&   Super User & \url{https://superuser.com/} \\
&   Stack Overflow & \url{https://stackoverflow.com} \\
& Quora & \url{https://www.quora.com/} \\
&   Make-Use-Of & \url{https://www.makeuseof.com/} \\
&   YouTube & \url{https://www.youtube.com/} \\
\hline
\multirow{18}{*}{Workflow} & UniPath Marketplace & \url{https://marketplace.uipath.com/} \\
&   sync.blue & \url{https://www.sync.blue/} \\
&   Device Tests & \url{https://devicetests.com/} \\
&   Make Tech Easier & \url{https://www.maketecheasier.com/} \\
&   Unix \& Linux Stack Exchange & \url{https://unix.stackexchange.com/} \\
&   Geeks for Geeks & \url{https://www.geeksforgeeks.org/} \\
&   I Love Free Software & \url{https://www.ilovefreesoftware.com/} \\
&   The Geek Diary & \url{https://www.thegeekdiary.com/} \\
&   Zyxware & \url{https://www.zyxware.com/} \\
&   GNOME Discourse & \url{https://discourse.gnome.org/} \\
&   It's FOSS & \url{https://itsfoss.com/} \\
&   Super User & \url{https://superuser.com/} \\
&   Stack Overflow & \url{https://stackoverflow.com} \\
&   LibreOffice Forum & \url{https://ask.libreoffice.org/} \\
&   ImpressExtractNotes & \url{https://github.com/danielrcollins1/ImpressExtractNotes} \\
&    Medium & \url{https://medium.com/} \\
&   YouTube & \url{https://www.youtube.com/} \\
&   Kelvin Smith Library & \url{https://case.edu/library/} \\
\end{longtable}

\subsection{Task Examples Collection}
\label{app:examples_collection}

Here we show the detailed statistics of \ours benchmark, including the main set on Ubuntu (\numexamples in total) and the analytic set on Windows (\numexampleswindows in total).
\begin{table}[htbp]
\centering
\caption{
Detailed statistics of \ours benchmark suite about examples number, average instruction tokens, infeasible instructions and integrated instructions.
}
\label{tab:d-env-stats}
\resizebox{\textwidth}{!}{
\begin{tabular}{@{}lccccccccccc@{}}
\toprule
 & OS & Calc & Impress & Writer & VLC & Thunderbird & Chrome & VSCode & GIMP & \textbf{Workflow} & \textbf{Overall} \\
\midrule
Examples & 24 & 47 & 47 & 23 & 17 & 15 & 46 & 23 & 26 & 101 & 369 \\
% \xspace\xspace\xspace\xspace Easy & & & & & & & & & & & \\
% \xspace\xspace\xspace\xspace Medium & & & & & & & & & & & \\
% \xspace\xspace\xspace\xspace Hard & & & & & & & & & & & \\
% \xspace\xspace\xspace\xspace Extra Hard & & & & & & & & & & & \\
Avg. Inst. Tokens & 22.38 & 33.30 & 25.19 & 35.30 & 35.82 & 34.07 & 22.07 & 20.78 & 16.23 & 51.24 & 33.36 \\
\#Infeasible & 5 & 1 & 0 & 1 & 3 & 1 & 3 & 5 & 10 & 1 & 30 \\
\#Integrated & 7 & 19 & 30 & 0 & 0 & 0 & 26 & 0 & 0 & 2 & 84 \\
% Avg. Action Steps & & & & & & & & & & & \\
\bottomrule
\end{tabular}
}
\end{table}

\begin{table}[htbp]
\centering
\caption{
Detailed statistics of Windows analytic set benchmark suite. This set contains no infeasible tasks and integrated tasks.
}
\label{tab:d-env-stats-windows}
\resizebox{0.5\textwidth}{!}{
\begin{tabular}{@{}lccccc@{}}
\toprule
 & Excel & Word & PPT & \textbf{Workflow} & \textbf{Overall} \\
\midrule
Examples & 11 & 9 & 7 & 16 & 43 \\
Avg. Inst. Tokens & 19.45 & 21.44 & 21.86 & 47.57 & 32.48\\
\bottomrule
\end{tabular}
}
\end{table}

We also visualize the intent distribution (We obtain sentence embeddings for instructions using OpenAI's embedding model, and then apply t-SNE to reduce the dimensionality to two dimensions for visualization.) and compare it with other benchmarks which also focus on the digital agent.
We randomly sample 300 examples from each dataset and randomly choose three different hyperparameters for t-SNE.
Visualization results are shown in Figure~\ref{fig:instructions_distribution_visualization}.
From the figure, we can observe that the semantic distribution of the instructions alone has reached the most comprehensive level. 
Additionally, our environment remains controllable and executable, offering a more reliable evaluation. 
It is also noticeable that the clustering centers of the other three are closely positioned, whereas the points in our distribution approaches are inconsistent with theirs, indicating that we can serve as a unique choice for a more comprehensive assessment of the capabilities of future intelligent agents.

\begin{figure}
    \centering
    \includegraphics[width=0.32\linewidth]{images/instructions_visualization/tsne_1.jpg}
    \includegraphics[width=0.32\linewidth]{images/instructions_visualization/tsne_2.jpg}
    \includegraphics[width=0.32\linewidth]{images/instructions_visualization/tsne_3.jpg}
    \caption{Comparison of instructions distribution. All datasets are sampled to 300 to make a fair comparison. The hyper-parameters of t-SNE are randomly sampled for each plot.}
    \label{fig:instructions_distribution_visualization}
\end{figure}


\subsection{Initial State Setup Details}
\label{app:initial_state_setup_details}
The setup of the initial state contains three stages:
1) Start emulator. 
The specified virtual machine is activated and automatically reverted to the corresponding snapshot, which records the initial system settings of the machine.
2) Prepare files (Optional). The file or software that specifies the initial state of the task to be executed is downloaded to the virtual machine and opened. 
The system is configured to first download the files to the host through a direct link and then upload them to the VM via a LAN connection. 
Specifically, some initial files are set up for OS-related tasks by manipulating the file system directly from the command line.
3) Execute reprocessing commands (Optional).
For tasks that require additional preprocessing, task-specific operations are executed after the completion of the first two phases. 
For example, taking the currently open LibreOffice Impress file to page five, clicking in the center of the screen to return to the main interface, \textit{etc}.
We provide convenient APIs to configure initial conditions and world settings, standardizing our tasks to make this process user-friendly and easily extendable.

\subsection{Evaluation Configuration Details}
\label{app:evaluation_configuration_details}
In this section, we will show details of preparations for the evaluation of the selected apps (LibreOffice -- Calc, Writer and Impress, Thunderbird, VLC Media Player, Chrome, VS Code, GIMP) and OS (Ubuntu and Windows).

\subsubsection{Ubuntu}

\paragraph{LibreOffice: Calc, Writer, and Impress}
LibreOffice is a popular open-source fully-featured office suite for Linux desktops. Our benchmark is built upon version 7.3.7.2, the 
version pre-installed in Ubuntu 22.04. Calc, Writer, and Impress are selected to build tasks on them. As the majority of tasks are to 
conduct a little revision to a file, we evaluate these tasks mainly by checking the final result file (in xlsx, docx, or pptx format).
The check can be done by comparing 
the result file with a golden reference, or inspecting some particular attributes or settings of the file, \textit{e.g.,}, page style,
freezing, and locale. Usually, the xlsx, docx, and pptx files are mainly accessed through 
\verb|openpyxl|\footnote{\url{https://openpyxl.readthedocs.io/en/stable/}},
\verb|python-docx|\footnote{\url{https://github.com/python-openxml/python-docx}}, and
\verb|python-pptx|\footnote{\url{https://github.com/scanny/python-pptx}}. For some properties not supported by the current libraries,
we also look them up directly via parsing the Office Open XML 
format\footnote{\url{https://learn.microsoft.com/en-us/office/open-xml/about-the-open-xml-sdk}}.

\paragraph{Thunderbird}
Thunderbird is a popular open-source fully-featured email client for Linux desktops. Version 115.6.0 of Thunderbird is pre-installed
in Ubuntu 22.04. We crafted an account profile to set up a feasible initial state. Evaluation for Thunderbird configurations
is mainly performed by reading various configurations or data files in the profile folder. An open-source reverse engineering
tool Firefox Decrypt\footnote{\url{https://github.com/unode/firefox_decrypt}} is adopted to decrypt the stored account 
information for evaluation. The involved account information is just for examples and contains no information about the real person.
Besides, there are tasks instructing to help to compose a new email. In these cases, the accessibility tree is leveraged to inspect 
the contents in the composing window before really sending it.


\paragraph{VLC Media Player}
VLC Media Player is a popular open-source cross-platform multimedia player and framework that plays most multimedia files. The evaluation for VLC Media Player is multifold, ranging from utilizing VLC HTTP interface\footnote{\url{https://wiki.videolan.org/Control_VLC_via_a_browser/}}, reading the VLC configuration file, comparing final result files, and leveraging accessibility tree to inspect the desired content.

\paragraph{Chrome}
Google Chrome is one of the most popular and powerful cross-platform web browsers developed on Google's free and open-source software project Chromium. 
The evaluation of Chrome is mainly based on the utilization of Playwright\footnote{\url{https://playwright.dev/}}, a browser automation library to control Chromium, Firefox, and WebKit with a single API. To connect Playwright running 
on host machine with Chrome running on virtual machine, port transferring tool 
socat\footnote{\url{http://www.dest-unreach.org/socat/},\url{https://linux.die.net/man/1/socat}} 
is leveraged. Additional information such as the HTML source codes of websites is also leveraged in the evaluation of some tasks.

\paragraph{VS Code}
VS Code is a popular open-source multi-functional cross-platform editor for source-code editing. 
The evaluation of VS Code tasks is primarily divided into two different categories. One subset of tasks is predominantly oriented towards file manipulation. 
In the context of these tasks, a comparative analysis is conducted between the resultant file and an anticipated reference gold file. 
Another subset of tasks is centered around how to utilize the intrinsic functionalities of the VS Code software itself, such as modifying color themes, initiating workspace sessions, and modifying settings. 
In these instances, it becomes important to extract relevant internal information and configurations from the VS Code environment.

In the course of this research, we principally leverage the capabilities offered by the VS Code Extension API\footnote{\url{https://code.visualstudio.com/api}} and information in the settings JSON file\footnote{\url{https://code.visualstudio.com/docs/getstarted/settings\#_settingsjson}} to obtain the requisite internal signal for the evaluation process. Our methodology involves the development of a custom VS Code extension, its installation within the VS Code software deployed on our virtual machine, and the subsequent invocation of the extension's command each time an evaluation is required, as well as checking whether the settings JSON has the correct value for a specific key.

\paragraph{GIMP}
GIMP is an open-source raster graphics editor used for image manipulation, editing, free-form drawing, format transcoding, and more specialized tasks. 
The evaluation for GIMP tasks is also mainly divided into two different categories, just like the VS Code evaluation. 
One type of task is mainly oriented to file operations.
In these tasks, the resulting files are compared and analyzed with the expected reference golden files, mainly relying on some powerful image processing libraries such as \verb|pillow|\footnote{\url{https://pypi.org/project/pillow/}}. 
Another category of tasks revolves around taking advantage of the inherent capabilities of the GIMP software itself. 
In these instances, we primarily read GIMP's configuration files to obtain internal information to evaluate the tasks.


\subsection{Windows}

\paragraph{Microsoft Office: Excel, Word, and PowerPoint} Microsoft Office is the most
popular office suite on Windows desktops. These three components share the same functions
with the corresponding LibreOffice components by and large. They are used to edit xlsx, docx,
and pptx files, respectively. Thus, the evaluation for LibreOffice tasks can be reused
for Microsoft Office tasks.

\paragraph{Thunderbird} Thunderbird is a cross-platform email client. Only the 
structure of profile folder on Windows is sightly
different from that on Linux. We thus revised the account profile and reuse it to 
set up the same initial state on Windows.

\paragraph{Chrome} Chrome is a cross-platform web browser. To evaluate tasks on 
Chrome, only the port transferring tool needs to be replaced with 
Ncat\footnote{\url{http://www.dest-unreach.org/socat/}}. Other configurations and the
evaluations can be shared with Linux-version tasks.

\subsection{More Task Examples}
\label{app:representitive_examples}
In this section, we curate a collection of examples from various app sets, each characterized by distinct operational logic and requiring different capabilities. 
These examples are carefully chosen to illustrate the diverse challenges and requirements encountered when interacting with different types of applications.

\begin{longtable}{m{1.5cm}m{4cm}m{6cm}m{2cm}}
\caption{More Example Showcase from Each Subset of Domains.} \label{tab:representative_cases} \\
\textbf{Related App(s)} & \textbf{Instruction(s)} & \textbf{Screenshot} & \textbf{Abilities Needed} \\
\hline
\endfirsthead

\multicolumn{4}{c}%
{{\bfseries \tablename\ \thetable{} -- continued from previous page}} \\
\hline
\textbf{Related App(s)} & \textbf{Task Instruction} & \textbf{Screenshot of Initial State} & \textbf{Abilities Needed} \\
\hline
\endhead

\multicolumn{4}{r}{{\textit{Continued on next page}}} \\ 
\endfoot

\endlastfoot

% Example row, repeat this block for each case
OS & \textit{I want to install Spotify on my current system. Could you please help me?} & \includegraphics[width=6cm, height=3.38cm]{images/cases/os_install_spotify.png} &  \texttt{knowledge of OS; omit distractions} \\
\hline
Calc &  \textit{Check the names in column ``Names with duplicates'' and put the unique ones in column ``Unique Names''. Keep the original order.} & \includegraphics[width=6cm, height=3.38cm]{images/cases/calc_unique_names.png} & \texttt{massive elements; knowledge tricks or reasoning over long actions} \\
\hline
Calc & \textit{I have a lookup table for the officers of each branch. Please, here is another table in which I need to fill with the officer names according the headoffice (i.e., the branch name). Help me to complete this.} & \includegraphics[width=6cm, height=3.38cm]{images/cases/calc_lookup.png} & \texttt{massive elements; knowledge of formulas and functions} \\
\hline
Calc & \textit{Given a partial calendar, please highlight all the weekends (Saturday \& Sunday) by setting the cell background as red (\#ff0000).}  & \includegraphics[width=6cm, height=3.38cm]{images/cases/calc_highlight_dates.png} & \texttt{massive elements; commonsense reasoning; software tricks} \\
\hline
Impress & \textit{I closed the slide pannel on the left and idk how to get it back please help} & \includegraphics[width=6cm, height=3.38cm]{images/cases/impress_close_left_panel.png} &  \texttt{software knowledge; imagine about UI layouts; overcome typos in instruction} \\
\hline
Impress & \textit{On it Whenever I launch a LibreOffice Impress, it uses both screens, one for current slide and next slide and another for actual presentation. What I want is to use only one monitor which shows presentation. I dont want the screen with Current slide and Next slide so that it can be used for other purposes. How should I achieve this?} & \includegraphics[width=6cm, height=3.38cm]{images/cases/impress_2_screens.png} &  \texttt{reason from unprofessional phenomenon expression} \\
\hline
Writer & \textit{Copy the screenshot 1.png from the desktop to where my cursor is located} & \includegraphics[width=6cm, height=3.38cm]{images/cases/writer_insert_picture.png} & \texttt{locate the position of cursor; switch from desktop and app} \\
\hline
Chrome &  \textit{Can you help me clean up my computer by getting rid of all the tracking things that Amazon might have saved? I want to make sure my browsing is private and those sites don't remember me.} & \includegraphics[width=6cm, height=3.38cm]{images/cases/chrome_clean_cookies.png} &  \texttt{understanding the unprofessional expression} \\
\hline
VLC & \textit{I am reading lecture note in PDF while a music video is running in VLC media player. But I find I need to switch to the player every time I need to pause/start.Could you help me change the setting to allow pausing the video using keyboard shortcut without minimizing the PDF reader? I want to focus on the lecture note and don't be disturbed by the app switching.} & \includegraphics[width=6cm, height=3.8cm]{images/cases/vlc_background_music.png} & \texttt{understanding the reference from unprofessional expression; software knowledge} \\
\hline
VLC & \textit{Hey, could you turn this video the right way up for me? And once it's flipped around, could you save it for me with the name `1984\_Apple.mp4' on the main screen where all my files are?} & \includegraphics[width=6cm, height=3.8cm]{images/cases/vlc_upside_down.png} &  \texttt{software knowledge; spatial judgment ability} \\
\hline
Thunderbird & \textit{Create a local folder called "Promotions" and create a filter to auto move the inbox emails whose subject contains “discount” to the new folder} & \includegraphics[width=6cm, height=3.38cm]{images/cases/thunderbird_promotion_file.png} & \texttt{software knowledge} \\
\hline
Thunderbird & \textit{Attach the my AWS bill to the email. The bill is stored at ~/aws-bill.pdf. Don't close it or send it. I haven't finish all the contents.} & \includegraphics[width=6cm, height=3.38cm]{images/cases/thunderbird_attach_aws.png} & \texttt{file management; extra requirement} \\
\hline
Thunderbird & \textit{I've got a bunch of email accounts in Thunderbird, and it's a hassle to check them one by one. Can you show me how to set up a unified inbox so I can see all my emails in one place?} & \includegraphics[width=6cm, height=3.38cm]{images/cases/thunderbird_unifiedbox.png} & \texttt{deep-hided feature, need to be explored even by human users; pop-up window} \\
\hline
VS Code & \textit{Please modify VS Code's settings to disable error reporting for Python missing imports.} & \includegraphics[width=6cm, height=3.38cm]{images/cases/vscode_disable_error_reporting.png} & \texttt{software knowledge to deal with settings; reasoning to understand the cause and solution of the error} \\
\hline
VS Code & \textit{Please help me install the autoDocstring extension in VS Code.} & \includegraphics[width=6cm, height=3.38cm]{images/cases/vscode_autostring.png} & \texttt{software knowledge to deal with Extensions; reasoning to search and install the extension successfully} \\
\hline
GIMP & \textit{Could you make the background of this image transparent for me?} & \includegraphics[width=6cm, height=3.38cm]{images/cases/gimp_remove_bg.png} & \texttt{precise and intricate operations} \\
\hline
GIMP & \textit{Help me choose the yellow triangle and position it at the center of my picture.} & \includegraphics[width=6cm, height=3.38cm]{images/cases/gimp_move_triangle.png} &  \texttt{spatial perception and reasoning, as well as precise control of actions} \\
\hline
Multiple {\fontsize{7pt}{10pt}\selectfont (VLC+GIMP)} & \textit{Could you help me create an Animated GIF from a video file using VLC and GIMP from the source of video ``src.mp4'', 5-second clip beginning at 00:03?} & \includegraphics[width=6cm, height=3.38cm]{images/cases/multi_apps_create_gif.png} & \texttt{software knowledge to undergo sophisticated processes and ability to process multi-step procedure successfully} \\
\hline
Multiple \thead{\fontsize{7pt}{10pt}\selectfont (ThunderBird+\\ \fontsize{7pt}{10pt}\selectfont Writer+Chrome)} & \textit{Help me export charts, graph or other images from docx files received in email ``Lecture Document'' in Notes folder and upload these png files to the figures/ folder in Google Drive for later use (use numbers to name them).} & \includegraphics[width=6cm, height=3.38cm]{images/cases/multi_apps_extract_from_file.png} & \texttt{ability to selectively export charts, graphs and images from docx file; software knowledge for google drive file upload} \\
\hline
Multiple {\fontsize{7pt}{10pt}\selectfont (Chrome+Calc)} & \textit{Could you help me extract data in the table from a new invoice uploaded to my Google Drive, then export it to a Libreoffice calc .xlsx file in the desktop?} & \includegraphics[width=6cm, height=3.38cm]{images/cases/multi_apps_extract_invoice.png} & \texttt{ability to do table data extraction; export data to .xlsx file} \\
\hline
% Multiple &  & \includegraphics[width=6cm, height=3.38cm]{images/environment/desktop_screenshot.png} & \\
% \hline
% Add more rows as needed
\end{longtable}



\section{Details of Baseline Methods}
\subsection{Hyper-Parameter of the Baseline Agents}
\label{app:prompting_details}
We utilize the versions of \texttt{gpt-3.5-turbo-16k}, \texttt{gpt-4-0125-preview}, and \texttt{gpt-4-vision-preview}, respectively for GPT results, need to be noted that result could be changed from time since it is close-sourced.
We also employ the \texttt{gemini-pro} and \texttt{gemini-pro-vision} versions for the Gemini models 
For all language models, we set the temperature parameter to 1.0, and top\_p to 0.9, and the maximum number of tokens for generation is set to 1500. 
We set the maximum steps of interaction to 15 and the maximum time limits to 30 minutes for all tasks since the agent could lead to a stuck environment under some unexpected cases.

\subsection{Prompt Details}

\subsubsection{Prompt for A11y Tree, Screenshot and their Combination Setting}

\begin{tcolorbox}
\begin{Verbatim}[breaklines=true]
You are an agent which follow my instruction and perform desktop computer tasks as instructed.
You have good knowledge of computer and good internet connection and assume your code will run on a computer for controlling the mouse and keyboard.
For each step, you will get an observation of an image, which is the screenshot of the computer screen and you will predict the action of the computer based on the image.

You are required to use `pyautogui` to perform the action grounded to the observation, but DONOT use the `pyautogui.locateCenterOnScreen` function to locate the element you want to operate with since we have no image of the element you want to operate with. DONOT USE `pyautogui.screenshot()` to make screenshot.
Return one line or multiple lines of python code to perform the action each time, be time efficient. When predicting multiple lines of code, make some small sleep like `time.sleep(0.5);` interval so that the machine could take; Each time you need to predict a complete code, no variables or function can be shared from history
You need to to specify the coordinates of by yourself based on your observation of current observation, but you should be careful to ensure that the coordinates are correct.
You ONLY need to return the code inside a code block, like this:
```python
# your code here
```
Specially, it is also allowed to return the following special code:
When you think you have to wait for some time, return ```WAIT```;
When you think the task can not be done, return ```FAIL```, don't easily say ```FAIL```, try your best to do the task;
When you think the task is done, return ```DONE```.

My computer's password is 'password', feel free to use it when you need sudo rights.
First give the current screenshot and previous things we did a short reflection, then RETURN ME THE CODE OR SPECIAL CODE I ASKED FOR. NEVER EVER RETURN ME ANYTHING ELSE.
\end{Verbatim}
\end{tcolorbox}

For a11y tree setting and a11y tree + screenshot setting, the prompts are basically the same, just replace the screenshot words with a11y tree words.

\subsubsection{Prompt for SoM Setting}
\begin{tcolorbox}
\begin{Verbatim}[breaklines=true]
You are an agent which follow my instruction and perform desktop computer tasks as instructed.
You have good knowledge of computer and good internet connection and assume your code will run on a computer for controlling the mouse and keyboard.
For each step, you will get an observation of the desktop by 1) a screenshot with interact-able elements marked with numerical tags; and 2) accessibility tree, which is based on AT-SPI library. And you will predict the action of the computer based on the image and text information.

You are required to use `pyautogui` to perform the action grounded to the observation, but DONOT use the `pyautogui.locateCenterOnScreen` function to locate the element you want to operate with since we have no image of the element you want to operate with. DONOT USE `pyautogui.screenshot()` to make screenshot.
You can replace x, y in the code with the tag of the element you want to operate with. such as:
```python
pyautogui.moveTo(tag_3)
pyautogui.click(tag_2)
pyautogui.dragTo(tag_1, button='left')
```
When you think you can directly output precise x and y coordinates or there is no tag on which you want to interact, you can also use them directly. 
But you should be careful to ensure that the coordinates are correct.
Return one line or multiple lines of python code to perform the action each time, be time efficient. When predicting multiple lines of code, make some small sleep like `time.sleep(0.5);` interval so that the machine could take; Each time you need to predict a complete code, no variables or function can be shared from history
You need to to specify the coordinates of by yourself based on your observation of current observation, but you should be careful to ensure that the coordinates are correct.
You ONLY need to return the code inside a code block, like this:
```python
# your code here
```
Specially, it is also allowed to return the following special code:
When you think you have to wait for some time, return ```WAIT```;
When you think the task can not be done, return ```FAIL```, don't easily say ```FAIL```, try your best to do the task;
When you think the task is done, return ```DONE```.

My computer's password is 'password', feel free to use it when you need sudo rights.
First give the current screenshot and previous things we did a short reflection, then RETURN ME THE CODE OR SPECIAL CODE I ASKED FOR. NEVER EVER RETURN ME ANYTHING ELSE.
\end{Verbatim}
\end{tcolorbox}

\subsection{Accessibility Tree Filtering}
\label{app:a11y_tree_handling}
Since the original tree is large (usually over 1 million tokens in XML format), we 
filter the accessibility tree nodes by their tags, visibilities, availabilities,
\textit{etc}. The concrete rules are illustrated in 
the following Table~\ref{table:a11y_tree_handling}.

\begin{longtable}{p{0.3\linewidth}p{0.3\linewidth}p{0.3\linewidth}}
\caption{Criteria for keeping a11y tree nodes on Ubuntu and Windows platforms}
\label{table:a11y_tree_handling}\\
\toprule
\textbf{Condition} & \textbf{Ubuntu} & \textbf{Windows} \\
\midrule
\endfirsthead
\toprule
\textbf{Condition} & \textbf{Ubuntu} & \textbf{Windows} \\
\midrule
\endhead
\bottomrule
\endfoot
Node Tags & document$*$, $*$item, $*$button, $*$heading, $*$label, $*$scrollbar, $*$searchbox, $*$textbox, $*$link, $*$tabelement, $*$textfield, $*$textarea, $*$menu,
alert, canvas, check-box, combo-box, entry, icon, image, paragraph, scroll-bar, section, slider, static, table-cell, terminal, text, netuiribbontab, start, trayclockwclass, traydummysearchcontrol, uiimage, uiproperty, uiribboncommandbar & Same as Ubuntu \\
\midrule
Showing & True & Not Applicable \\
\midrule
Visible & True & True \\
\midrule
Enabled or Editable or Expandable or Checkable & True & True \\
\midrule
Has Name or Text or Image & True & True \\
\midrule
Position &  $>= 0$ &    $>= 0$ \\
\midrule
Size &  $> 0$ &    $> 0$ \\
\end{longtable}


\subsection{Set-of-Mark Implementation Details}
\label{app:som_details}
Our methodology involves an initial transformation of the original observational image acquired from our environment into the standardized \textit{SoM} format and be putted into VLMs together with a table of the marks with metadata information such as tags and names.
This format consists of bounding boxes that bound the sub-images of interest, each associated with a corresponding integer mark. 
Notably, our approach incorporates the utilization of the accessibility tree to identify the bounding boxes associated with all clickable buttons within the current image observation, instead of using segmentation models like the original \textit{SoM} paper.
Some examples of screenshots after applying SoM are shown in Figure~\ref{fig:som_exampels}.
We can observe the emergence of some low-quality, unmodelable tasks, and even misleading bounding boxes, depending on the level of support from the software ecosystem.
This could be another reason for the poor performance of SoM. Future improvements could be made in this regard.

\begin{figure}[htbp]
    \centering
    \includegraphics[width=0.32\linewidth]{images/som_examples/0c35396a-e6cd-44da-9b48-0d53d20d604b.png}
    \includegraphics[width=0.32\linewidth]{images/som_examples/9dc4bd8c-3643-40e5-9474-1b4aca6d9d9b.png}
    \includegraphics[width=0.32\linewidth]{images/som_examples/9dc917c1-dc12-4738-8058-a271bda008de.png}
    \includegraphics[width=0.32\linewidth]{images/som_examples/9df3e8c8-cdfc-48c9-83df-f80b1a4f2600.png}
    \includegraphics[width=0.32\linewidth]{images/som_examples/9df8d5e4-f7b4-4bdf-b39d-a31a2e7eeae8.png}
    \includegraphics[width=0.32\linewidth]{images/som_examples/c8731e48-a042-4d9e-be23-172c92637597.png}
    \caption{
    Showcase of example screenshots marked by SoM across different app GUI.
    }
    \label{fig:som_exampels}
\end{figure}

\subsection{Full Results of Baseline Methods}
\label{app:full_results}
Here we show the break-down results of baseline methods from different LLMs and VLMs for follow-up reference.

We have also compiled the distribution of steps taken by the GPT-4V model under our four settings: Accessibility Tree (A11y Tree), Screenshot, Screenshot combined with Accessibility Tree (Screenshot+A11y Tree), and Set-of-Mark. 
This data (as shown in Fig.~\ref{fig:steps_distribution}) provides potential clues for future work. 
Overall, there are observable differences in how many steps the agent chooses to execute and when it decides to terminate under different settings. 
More detailed control and analysis of these behaviors can be explored in subsequent research efforts.

\begin{figure}[htbp]
    \centering
    \includegraphics[width=0.8\linewidth]{images/agent_steps_hist.pdf}
    \caption{
    Distribution of steps taken by the GPT-4V based agents across four different settings.
    }
    \label{fig:steps_distribution}
\end{figure}

\begin{table*}[ht]
\centering
\caption{
Detailed success rates of baseline LLM and VLM agents on \ours, divided by apps (domains): OS, LibreOffice Calc, LibreOffice Impress, LibreOffice Writer, Chrome, VLC Player, Thunderbird, VS Code, GIMP and Workflow which is comprehensive with multiple apps, for gaining insights from interfaces and operation logics.
}
\label{tab:exp_baseline_full}
\small
\scalebox{1}{
\renewcommand{\arraystretch}{1.1} % Increase the row height
\begin{tabularx}{\linewidth}{>{\centering\arraybackslash}p{1cm}>{\centering\arraybackslash}p{2cm}>{\centering\arraybackslash}X>{\centering\arraybackslash}X>{\centering\arraybackslash}X>{\centering\arraybackslash}X>{\centering\arraybackslash}X>{\centering\arraybackslash}X>{\centering\arraybackslash}X>{\centering\arraybackslash}X>{\centering\arraybackslash}X>{\centering\arraybackslash}X}
\toprule
\multirow{2}{*}{\textbf{Inputs}} & \multirow{2}{*}{\textbf{Model}} & \multicolumn{9}{c}{\textbf{Success Rate (↑)}} \\ 
\cline{3-12}
~ & ~ & OS & Calc &Impress & Writer & VLC & TB & Chrome & VSC & GIMP & Workflow \\
\midrule
A11y & Mixtral-8x7B  & 12.50 & 0.00 & 0.39 & 4.34  & 10.22 &6.67 & 2.17 & 8.69 & 3.85 & 0.10 \\
& GPT-3.5   & 4.17 & 2.13 & 6.77 & 4.35 & 6.53 & 0.00 & 2.17 & 0.00 & 0.00 & 1.62 \\
& Gemini-Pro   & 4.17 & 0.00 & 2.13 & 4.35 & 12.41 & 0.00 & 2.17 & 0.00 & 7.69 & 0.63 \\
% & GPT-4-1106 &  & &  & & & \\
& GPT-4   & 20.83 & 0.00 & 6.77 & 4.35 & 23.53 & 26.67 & 26.09 & 30.43 & 23.08 & 2.97 \\
% & Gemini-Pro-1.5 &  & &  & & & \\
& Gemini-Pro-1.5   & 12.50 & 2.13 & 2.13 & 4.35 & 6.53 & 0.00 & 10.87 & 8.70 & 0.00 & 3.60 \\
% & Llama-3-70B &  & &  & & & \\
& Llama-3-70B & 4.17 & 0.00 & 0.39 & 8.70 & 6.53 & 0.00 & 2.17 & 0.00 & 0.00 & 0.63 \\
% & GPT-4o &  & &  & & & \\
& GPT-4o   & 20.83 & 6.38 & 6.77 & 8.69 & 12.41 & 20.00 & 17.39 & 21.74 & 11.54 & 7.56 \\
% & Qwen-Max &  & &  & & & \\
& Qwen-Max   & 29.17 & 0.00 & 2.52 & 13.04 & 8.95 & 0.00 & 10.87 & 8.70 & 11.54 & 2.61 \\
\midrule
Screen & CogAgent    & 4.17 & 0.00 & 0.00 & 4.34 & 6.53 &0.00 & 2.17 & 0.00 & 0.00 & 0.00 \\
& Gemini-ProV  & 8.33 & 0.00 & 6.77 & 4.35 & 12.41 & 0.00 & 6.52 & 8.70 & 23.08 & 2.08 \\
& GPT-4V    & 12.50 & 0.00 & 2.52 & 4.35 & 18.34 & 0.00 & 6.52 & 0.00 & 7.69 & 6.04 \\
& Claude-3-Opus  & 4.17 & 0.00 & 2.52 & 4.34 & 6.53 &0.00 & 2.17 & 0.00 & 3.84 & 2.61 \\
% & Gemini-Pro-1.5 &  & &  & & & \\
& Gemini-Pro-1.5   & 12.50 & 0.00 & 13.16 & 8.70 & 6.53 & 0.00 & 2.17 & 0.00 & 11.54 & 3.60 \\
% & GPT-4o &  & &  & & & \\
& GPT-4o   & 8.33 & 0.00 & 6.77 & 4.35 & 16.10 & 0.00 & 4.35 & 4.35 & 3.85 & 5.58 \\
\midrule
Screen & CogAgent  & 4.17 & 2.17 & 0.00 & 4.35 & 6.53 &0.00 & 2.17 & 0.00 & 0.00 & 0.10 \\
+ A11y & Gemini-ProV  & 4.17 & 2.13 & 6.77 & 4.35 & 18.30 & 0.00 & 4.35 & 0.00 & 0.00 & 1.52 \\
& GPT-4V  & 16.67 & 0.00 & 6.77 & 21.73 & 24.18 & 33.33 & 21.74 & 21.74 & 15.38 & 4.59 \\
& Claude-3-Opus  & 12.50 & 2.13 & 4.65 & 4.34 & 18.30 &0.00 & 2.17 & 8.69 & 7.69 & 0.99 \\
% & Gemini-Pro-1.5 &  & &  & & & \\
& Gemini-Pro-1.5   & 12.50 & 0.00 & 4.65 & 8.70 & 12.41 & 0.00 & 8.70 & 4.35 & 11.54 & 1.56 \\
% & GPT-4o &  & &  & & & \\
& GPT-4o   & 41.67 & 4.26 & 6.81 & 8.70 & 9.50 & 6.67 & 15.22 & 30.43 & 0.00 & 7.46 \\
\midrule
SoM & CogAgent  & 4.17 & 2.17 & 0.00 & 4.34 & 6.53 & 0.00 & 2.17 & 0.00 & 0.00 & 0.00 \\
& Gemini-ProV  & 4.17 & 0.00 & 0.39 & 4.34 & 6.53 & 0.00 & 0.00 & 0.00 & 0.00 & 0.63 \\
& GPT-4V  & 8.33 & 8.51 & 6.38 & 13.04 & 10.66 & 33.33 & 23.91 & 17.39 & 11.54 & 6.57 \\
& Claude-3-Opus  & 12.50 & 2.13 & 0.39 & 8.70 & 6.53 &13.33 & 17.39 & 0.00 & 11.54 & 4.49 \\
% & Gemini-Pro-1.5 &  & &  & & & \\
& Gemini-Pro-1.5   & 16.67 & 0.00 & 10.64 & 4.35 & 24.18 & 13.33 & 8.70 & 4.35 & 15.38 & 3.60 \\
% & GPT-4o &  & &  & & & \\
& GPT-4o   & 20.83 &0.00 & 6.77 & 4.35 & 6.53 & 0.00 & 4.35 & 4.35 & 0.00 & 3.60 \\
\midrule
\rowcolor{gray!20}
\multicolumn{2}{c}{Human Performance}  & 75.00 & 61.70 & 80.85 & 73.91 & 70.59 & 46.67 & 78.26 & 73.91 & 73.08 & 73.27 \\
% \rowcolor{gray!20}
% \multicolumn{2}{c}{Human (Expert)}  &  &  &  & \\
\bottomrule
\end{tabularx}
}
\end{table*}




% \section{Analysis on Different Input Settings}
% Here we provide a specific code execution process for each setting (Screenshot, A11y Tree, Set-of-Mark) in our main experiment and try to get insights for further improvement.

% \subsection{Effect of Screenshot}

% \begin{figure}[htbp]
%     \centering
%     \includegraphics[width=0.8\textwidth]{images/qualititive_analysis/easy_done/step_1_20240130@033529.png}
%     \caption{gimp/f723c744-e62c-4ae6-98d1-750d3cd7d79d}
%     \label{fig:error_case_gimpf723}
% \end{figure}


% \footnotesize
% \begin{Verbatim}
% # ACTION 1
% import pyautogui
% import time

% # Move to 'Colors' in the menu bar
% pyautogui.moveTo(400, 40, duration=1) 
% pyautogui.click()
% time.sleep(0.5)

% # Move to 'Brightness-Contrast'
% pyautogui.moveTo(400, 140, duration=1) 
% pyautogui.click()


% # ACTION 2
% import pyautogui
% import time

% # Move to 'Colors' in the menu bar
% pyautogui.moveTo(250, 30, duration=1)
% pyautogui.click()
% time.sleep(0.5)

% # Move to 'Brightness-Contrast'
% pyautogui.moveRel(0, 180, duration=1) 
% # Assuming 'Brightness-Contrast' is at this relative position from 'Colors'
% pyautogui.click()


% # ACTION 3
% import pyautogui
% import time

% # Move to 'Colors' in the menu bar, and click to open the dropdown
% pyautogui.moveTo(600, 25, duration=1) 
% # Y-coordinate adjusted to likely position of the 'Colors' menu
% pyautogui.click()
% time.sleep(0.5)

% # Move to 'Brightness-Contrast' in the dropdown and click
% pyautogui.moveRel(0, 160, duration=1)  
% # Y-coordinate adjusted to likely position of the 'Brightness-Contrast' option
% pyautogui.click()
% \end{Verbatim}




% \subsection{Effect of A11y Tree}

% \begin{figure}[htbp]
%     \centering
%     \includegraphics[width=0.8\textwidth]{images/qualititive_analysis/tree_info_benefits/step_1_20240130@172258.png}
%     \caption{os/7688b85f-87a4-4e4a-b2f8-f3d6c3f29b82}
%     \label{fig:error_case_os7688}
% \end{figure}


% \small
% \begin{Verbatim}
% todo
% \end{Verbatim}

% \subsection{Effect of Screenshot}

% \begin{figure}[htbp]
%     \centering
%     \includegraphics[width=0.8\textwidth]{images/qualititive_analysis/effect_of_screenshot/start.png}
%     \includegraphics[width=0.8\textwidth]{images/qualititive_analysis/effect_of_screenshot/right_next_step.png}
%     \caption{thunderbird/94760984-3ff5-41ee-8347-cf1af709fea0}
%     \label{fig:error_case_effect_of_screenshot}
% \end{figure}

% \small
% \begin{Verbatim}
% # Wrong action
% import pyautogui
% import time

% # Move the mouse to the "Install Theme" button for the "Dark Fox" theme and click it
% pyautogui.moveTo(945 + 136/2, 2810 + 36/2, duration=1)  # Aim for the center of the button
% time.sleep(0.5)
% pyautogui.click()
% \end{Verbatim}

% \subsection{Effect of Set-of-Mark}

% \begin{figure}[htbp]
%     \centering
%     \includegraphics[width=0.8\textwidth]{images/qualititive_analysis/som/5b6e243e-115d-436d-bde5-cbe739a1a800.png}
%     \caption{Tagged screenshot after applying Set-of-Mark.}
%     \label{fig:error_case_som_calc}
% \end{figure}


\section{Examples of Qualitative Analysis}
\label{appendix:qualitative_analysis}

Here we give the specific code execution process for the examples in the qualitative analysis in \S\ref{qualitative_analysis} and add more error cause analysis, trying to get insights for further improvement.

\subsection{Success and Failure Cases}

\begin{tcolorbox}[title={Success Task: \textit{I downloaded an episode of Friends to practice listening, but I don't know how to remove the subtitles. Please help me remove the subtitles from the video and export it as "subtitles.srt" and store it in the same directory as the video.}}]

\small
\begin{Verbatim}[commandchars=\\\{\}]
\textcolor{purple}{# Action 1}
import pyautogui
import time
\textcolor{blue}{# Move the mouse to the top bar of the media player window to access the menu}
pyautogui.moveTo(800, 10)  
\textcolor{blue}{# Assuming the menu is at the top of the screen}
pyautogui.click()
time.sleep(0.5)

\textcolor{purple}{# Action 2}
\textcolor{blue}{# Move the mouse to the Activities corner to open the GNOME Activities Overview}
pyautogui.moveTo(10, 10)  
\textcolor{blue}{# Assuming the Activities corner is at the top-left}
pyautogui.click()
\textcolor{blue}{# Type `terminal' to search for the Terminal application}
pyautogui.write(`terminal')
time.sleep(0.5)
\textcolor{blue}{# Press Enter to open the Terminal application}
pyautogui.press(`enter')

\textcolor{purple}{# Action 3}
\textcolor{blue}{# Focus the Terminal window and type the command to extract subtitles}
pyautogui.click(640, 360)  
\textcolor{blue}{# Clicking in the center of the screen to focus the Terminal}
time.sleep(0.5)
pyautogui.write(`ffmpeg -i video.mp4 -map 0:s:0 subtitles.srt')
pyautogui.press(`enter')

\textcolor{purple}{# Action 4}
\textcolor{blue}{# Focus the Terminal window and type the command}
\textcolor{blue}{ to create a new video without subtitles}
pyautogui.click(1100, 350)  
\textcolor{blue}{# Clicking inside the Terminal window to focus it}
pyautogui.write(`ffmpeg -i video.mp4 -c copy -sn no_subtitles_video.mp4')

\textcolor{purple}{Done!}
\end{Verbatim}
\end{tcolorbox}


\begin{tcolorbox}[title={Failure Task: \textit{Help me center align the heading in LibreOffice.}}]

\small
\begin{Verbatim}[commandchars=\\\{\}]
\textcolor{purple}{# Action 1}
import pyautogui
import time
\textcolor{blue}{# Move the mouse to the LibreOffice Writer document icon and double-click}
pyautogui.moveTo(1100, 720, duration=1)
pyautogui.doubleClick()

\textcolor{purple}{# Action 2}
\textcolor{blue}{# Move the mouse to the beginning of the heading text}
\textcolor{blue}{and click and drag to select the heading}
pyautogui.moveTo(200, 230, duration=1)
pyautogui.mouseDown()
pyautogui.moveTo(800, 230, duration=1)
pyautogui.mouseUp()
\textcolor{blue}{# Move the mouse to the center align button in the toolbar and click}
pyautogui.moveTo(400, 150, duration=1)
pyautogui.click()

\textcolor{purple}{# Action 3}
\textcolor{blue}{# Assuming the text is still selected, }
\textcolor{blue}{move the mouse to the center align button in the toolbar and click}
pyautogui.moveTo(400, 150, duration=1)
pyautogui.click()

\textcolor{purple}{# Action 4}
\textcolor{blue}{# Move the mouse to the center align button in the toolbar and click}
pyautogui.moveTo(340, 80, duration=1)
pyautogui.click()

\textcolor{purple}{Failed.}
\end{Verbatim}
\end{tcolorbox}

\begin{figure}[htbp]
    \centering
    \includegraphics[width=\linewidth]{images/qualititive_analysis/model_lack_knowledge.pdf}
    \caption{
    The figure shows examples of usage errors in three types of software. The first row is for GIMP, the second row is for LibreOffice Calc, and the third row is for Google Chrome. Detailed code execution processes are provided below.
    }
    \label{fig:model_lack_knowledge}
\end{figure}

\begin{tcolorbox}[title={GIMP Task: \textit{Could you tone down the brightness of my photo?}}]

\begin{Verbatim}[commandchars=\\\{\}]
\textcolor{purple}{# Action 1}
import pyautogui
import time
\textcolor{blue}{# Coordinates for the image icon on the desktop}
\textcolor{blue}{# Adding half of the width and height to get the center}
image_icon_x, image_icon_y = 1793 + 60, 920 + 17  
\textcolor{blue}{# Double-click on the image icon}
pyautogui.doubleClick(image_icon_x, image_icon_y)

\textcolor{purple}{# Action 2}
\textcolor{blue}{# Coordinates for the `Colors' menu}
\textcolor{blue}{# Adding half of the width and height to get the center}
colors_menu_x, colors_menu_y = 344 + 28, 64 + 12  
\textcolor{blue}{# Click on the `Colors' menu}
pyautogui.click(colors_menu_x, colors_menu_y)

\textcolor{purple}{# Action 3}
\textcolor{blue}{# Coordinates for the `Brightness-Contrast...' option in the `Colors' dropdown}
\textcolor{blue}{# Assuming each menu item is approximately 28 pixels apart}
brightness_contrast_option_x, brightness_contrast_option_y 
= 344 + 28, 64 + 12 + 28 * 6  
\textcolor{blue}{# Click on the `Brightness-Contrast...' option}
pyautogui.click(brightness_contrast_option_x, brightness_contrast_option_y)

\textcolor{purple}{# Action 4}
\textcolor{blue}{# Coordinates for the `Cancel' button in the `Shadows-Highlights' dialog}
\textcolor{blue}{# Adding half the width and height to get the center}
cancel_button_x, cancel_button_y = 375 + 42, 625 + 16  
\textcolor{blue}{# Click the `Cancel' button}
pyautogui.click(cancel_button_x, cancel_button_y)
\textcolor{blue}{# Coordinates for the `Brightness-Contrast...' option in the `Colors' dropdown}
\textcolor{blue}{# Adjusted index for Brightness-Contrast option}
brightness_contrast_option_x, brightness_contrast_option_y 
= 344 + 28, 64 + 12 + 28 * 5  
pyautogui.click(brightness_contrast_option_x, brightness_contrast_option_y)

\textcolor{purple}{# Action 5}
\textcolor{blue}{# Coordinates for the `Cancel' button in the `Exposure' dialog}
\textcolor{blue}{# Adding half the width and height to get the center}
cancel_button_x, cancel_button_y 
= 375 + 42, 391 + 16  
pyautogui.click(cancel_button_x, cancel_button_y)
\textcolor{blue}{# Coordinates for the `Brightness-Contrast...' option in the `Colors' dropdown}
\textcolor{blue}{# Adjusted index for Brightness-Contrast option}
brightness_contrast_option_x, brightness_contrast_option_y 
= 344 + 28, 64 + 12 + 28 $*$ 6  
pyautogui.click(brightness_contrast_option_x, brightness_contrast_option_y)

\textcolor{purple}{# Action 6}
\textcolor{blue}{# Coordinates for the `Cancel' button in the `Color Balance' dialog}
\textcolor{blue}{# Adding half the width and height to get the center}
cancel_button_x, cancel_button_y 
= 375 + 42, 588 + 16  
pyautogui.click(cancel_button_x, cancel_button_y)
\textcolor{blue}{# Since the `Brightness-Contrast...' option was not correctly opened before,}
\textcolor{blue}{we will use an estimated position based on the dropdown menu.}
\textcolor{blue}{# Assuming `Brightness-Contrast...' is the second option in the `Colors' menu.}
brightness_contrast_option_x, brightness_contrast_option_y 
= colors_menu_x, colors_menu_y + 28 * 2
pyautogui.click(brightness_contrast_option_x, brightness_contrast_option_y)

\textcolor{purple}{# Action 7}
\textcolor{blue}{# Coordinates for the `Cancel' button in the `Color Temperature' dialog}
\textcolor{blue}{# Adding half the width and height to get the center}
cancel_button_x, cancel_button_y = 375 + 42, 403 + 16  
pyautogui.click(cancel_button_x, cancel_button_y)
\textcolor{blue}{# Since the `Brightness-Contrast...' option was not correctly opened before, }
\textcolor{blue}{we will use an estimated position based on the dropdown menu.}
\textcolor{blue}{# Assuming `Brightness-Contrast...' is the second option in the `Colors' menu.}
brightness_contrast_option_x, brightness_contrast_option_y 
= colors_menu_x, colors_menu_y + 28 * 2
\textcolor{blue}{# Click on the `Brightness-Contrast...' option}
pyautogui.click(brightness_contrast_option_x, brightness_contrast_option_y)

\textcolor{blue}{# ... More Wrong Duplicate Actions; Can not find the right coordination.}

\textcolor{purple}{# Failed}
\end{Verbatim}
\end{tcolorbox}

\begin{tcolorbox}[title={Libreoffice calc Task: \textit{I now want to count the meeting cities of the three machine learning conferences in the past ten years from 2013 to 2019(including 2013 and 2019). I have listed the names and years of the conferences in excel. Please fill in the vacant locations.}}]


\begin{Verbatim}[commandchars=\\\{\}]
\textcolor{purple}{# Action 1}
\textcolor{blue}{# Move to cell C2 which is the first empty cell under `City' for ICLR 2013}
pyautogui.click(319, 222)  
\textcolor{blue}{# Coordinates of cell C2}
time.sleep(0.5)
\textcolor{blue}{# Type the city name for ICLR 2013}
pyautogui.write(`Scottsdale')
\textcolor{blue}{# Press Enter to move to the next cell}
pyautogui.press('enter')

\textcolor{purple}{# Action 2}
\textcolor{blue}{# Move to cell C3 which is the city cell for ICML 2013}
pyautogui.click(319, 268)  
\textcolor{blue}{# Coordinates of cell C3}
\textcolor{blue}{# Type the city name for ICML 2013}
pyautogui.write(`Atlanta')
\textcolor{blue}{# Press Enter to move to the next cell}
pyautogui.press('enter')

\textcolor{purple}{# Action 3}
\textcolor{blue}{# Move to cell C4 which is the city cell for NeurIPS 2013}
pyautogui.click(319, 289)  
\textcolor{blue}{# Coordinates of cell C4}
\textcolor{blue}{# Type the city name for NeurIPS 2013}
pyautogui.write(`Lake Tahoe')
\textcolor{blue}{# Press Enter to move to the next cell}
pyautogui.press('enter')"

\textcolor{blue}{... More copy-and-paste operations}

\textcolor{purple}{# Action 15}
\textcolor{blue}{# Move to cell C16 which is the city cell for NeurIPS 2017}
pyautogui.click(319, 541)  
\textcolor{blue}{# Coordinates of cell C16}
\textcolor{blue}{# Type the city name for NeurIPS 2017}
pyautogui.write(`Long Beach')
\textcolor{blue}{# Press Enter to move to the next cell}
pyautogui.press('enter')

\textcolor{purple}{Failed.}
\end{Verbatim}
\end{tcolorbox}


\begin{tcolorbox}[title={Chrome Task: \textit{Can you make Bing the main search thingy when I look stuff up on the internet?}}]

\begin{Verbatim}[commandchars=\\\{\}]
\textcolor{purple}{# Action 1}
pyautogui.click(x=1280, y=50)  
\textcolor{blue}{# Approximate coordinates for the `More' button}
pyautogui.move(0, 150)  
\textcolor{blue}{# Moving the cursor down to `Settings' in the drop-down menu}
pyautogui.click()  
\textcolor{blue}{# Clicking on `Settings'}

\textcolor{purple}{# Action 2}
pyautogui.click(tag_31)  
\textcolor{blue}{# Click on the `Finish update' button to}
\textcolor{blue}{approximate the location of `More' button}
pyautogui.move(-50, 0)  
\textcolor{blue}{# Move left from the `Finish update' button to approximate the `More' button}
pyautogui.click()  
\textcolor{blue}{# Click on the `More' button}
time.sleep(0.5)
pyautogui.move(0, 150)  
\textcolor{blue}{# Move down to the `Settings' option in the drop-down menu}
pyautogui.click()  
\textcolor{blue}{# Click on `Settings'}

\textcolor{purple}{# Action 3}
\textcolor{blue}{# Click on the `Customise Chrome' button which should be the `More' button}
time.sleep(1)  
\textcolor{blue}{# Wait for the dropdown menu to appear}
\textcolor{blue}{# The `Settings' option is usually located towards the bottom of}
\textcolor{blue}{the dropdown menu, so I'll move the cursor down to reach it}
pyautogui.move(0, 200)  
\textcolor{blue}{# Move down to the `Settings' option in the drop-down menu}
pyautogui.click()  
\textcolor{blue}{# Click on `Settings'}

\textcolor{purple}{# Action 4}
pyautogui.click(tag_47)  
\textcolor{blue}{# Click on the `Close' button to close the `Customise Chrome' menu}
pyautogui.click(tag_25)  
\textcolor{blue}{# Click on the `New Tab' button to open a new tab}
time.sleep(1)  
\textcolor{blue}{# Wait for the new tab to open}
pyautogui.typewrite(`chrome://settings/searchEngines')  
\textcolor{blue}{# Type the direct URL to the search engine settings page}
pyautogui.press(`enter')

\textcolor{purple}{# Action 5}
\textcolor{blue}{# Simulate pressing the `Tab' key multiple times to navigate through the options}
for _ in range(10):  
\textcolor{blue}{# Assuming it might take around 10 tabs to reach the search engine list}    
    pyautogui.press(`tab')    
    time.sleep(0.5)  
\textcolor{blue}{# Wait for the focus to move to the next element}
\textcolor{blue}{# Simulate pressing the `Enter' key to select an option,}
\textcolor{blue}{which would be done once we reach the Bing option}
pyautogui.press('enter')

\textcolor{purple}{# Action 6}
pyautogui.click(tag_42)  
\textcolor{blue}{# Click on the `Cancel' button to exit the edit mode for the '360' search engine}

\textcolor{purple}{Failed.}
\end{Verbatim}
\end{tcolorbox}

\newpage

\subsection{Common Errors by GPT-4V Agents }
\label{appendix:common_errors}

In the following two figures ( Figure \ref{fig:cookie_block} and Figure \ref{fig:wrong_click}), we take the web page tasks as examples to show the two most common types of errors in GPT4-V:
\textit{mouse click inaccuracies} and \textit{inadequate handling of environmental noise}.
\begin{figure}[htbp]
    \centering
    \includegraphics[width=0.8\linewidth]{images/qualititive_analysis/wrong_click.pdf}
    \caption{
    The error shown in the figure is due to mouse click inaccuracies. The agent was supposed to click on the product category images of the Nike official store and follow the instructions to search for women's jerseys priced over 60 dollars. However, due to a coordinate calculation error, it mistakenly clicked on the favorite button, creating a significant discrepancy between the current state and the target state. Under these circumstances, the agent is unable to backtrack to the previous state and start over.
    }
    \label{fig:wrong_click}
\end{figure}


\begin{figure}[htbp]
    \centering
    \includegraphics[width=0.8\linewidth]{images/qualititive_analysis/cookie_block.pdf}
    \caption{
    The error shown in the figure is due to inadequate handling of environmental noise. When the agent encounters web pages with cookies or pop-up ads, its unfamiliarity with the structure of such pages leads to attempts to interact directly with the web page without bypassing the cookies or pop-ups. Since the pop-ups or cookie requests are not closed, all actions are ineffective.
    }
    \label{fig:cookie_block}
\end{figure}

% \small

% \textbf{\textit{mouse click inaccuracies}}

% \begin{Verbatim}[commandchars=\\\{\}]
% \end{Verbatim}

% \textbf{\textit{repetitive steps}}

% \begin{Verbatim}[commandchars=\\\{\}]
% \end{Verbatim}

% \textbf{\textit{inadequate handling of environmental noise}}


\newpage

\subsection{Discrepancies in Task Difficulty between Agent and Human}
\label{appendix:human_vs_agent}

\begin{figure}[htbp]
    \centering
    \includegraphics[width=\linewidth]{images/qualititive_analysis/human_agent_supplement.pdf}
    \caption{
    Supplementary examples of tasks performed by humans and agents.
    }
    \label{fig:human_agent_supplement}
\end{figure}



\begin{tcolorbox}[title={Human > Agent Task: \textit{erase all the highlighted marks in this document}}]

\small
\begin{Verbatim}[commandchars=\\\{\}]
\textcolor{purple}{# Action 1}
import pyautogui
import time
\textcolor{blue}{# Move to the `Character Highlighting Color' button and click}
pyautogui.moveTo(1020, 122)
pyautogui.click()

\textcolor{purple}{# Action 2}
\textcolor{blue}{# Assuming the `No Color' option is the first option in the dropdown}
\textcolor{blue}{# Move the mouse slightly below the `Character Highlighting Color' button}
\textcolor{blue}{to click on `No Color'}
pyautogui.moveTo(1020, 180)
time.sleep(0.5)
pyautogui.click()

\textcolor{purple}{# Action 3}
\textcolor{blue}{# Move to the `Save' button and click}
pyautogui.moveTo(184, 108) 
time.sleep(0.5)
pyautogui.click()

\textcolor{purple}{# Failed.}
\end{Verbatim}
\end{tcolorbox}

\begin{tcolorbox}[title={Human > Agent Task: \textit{On page 2, add a note same to the title into the slide. Make the font of title bold.}}]

\begin{Verbatim}[commandchars=\\\{\}]
\textcolor{purple}{# Action 1}
\textcolor{blue}{# Click on the vertical scrollbar to scroll down and reveal more slides}
pyautogui.click(tag_43)  

\textcolor{purple}{# Action 2}
\textcolor{blue}{# Click on the second slide thumbnail to select it}
pyautogui.click(tag_2) 

\textcolor{purple}{# Action 3}
\textcolor{blue}{# Click on the text box to add text}
pyautogui.click(tag_42)  
pyautogui.typewrite(`Team Members', interval=0.1)  
\textcolor{blue}{# Type the title}
pyautogui.hotkey(`ctrl', `a')  
\textcolor{blue}{# Select all text}
pyautogui.hotkey(`ctrl', `b')  
\textcolor{blue}{# Make the text bold}

\textcolor{purple}{# Action 4}
\textcolor{blue}{# Click on the Bold button to apply bold formatting}
pyautogui.click(tag_68)  

\textcolor{purple}{Failed.}
\end{Verbatim}
\end{tcolorbox}


\begin{tcolorbox}[title={Agent > Human Task: \textit{use GIMP to cut out the 2s to 4s part of a video}}]

\small
\begin{Verbatim}[commandchars=\\\{\}]
\textcolor{purple}{# Action 1}
import pyautogui
import time
\textcolor{blue}{# Move the mouse to the assumed terminal icon position}
\textcolor{blue}{at the bottom of the dock and click to open it}
\textcolor{blue}{# Assuming the terminal icon is at the bottom of the dock}
pyautogui.moveTo(70, 1350, duration=1)  
pyautogui.click()
\textcolor{blue}{# Type the command to use ffmpeg to trim the video}
ffmpeg_command = `ffmpeg -ss 00:00:02 -to 00:00:04 -i /home/user/fullvideo.mp4 
-c copy /home/user/trimmedvideo.mp4'
pyautogui.write(ffmpeg_command, interval=0.05)
pyautogui.press('enter')

\textcolor{purple}{# Action 2}
\textcolor{blue}{# Close the search dialog}
pyautogui.press('esc')
\textcolor{blue}{# Open the terminal using a keyboard shortcut (Ctrl+Alt+T is a common shortcut}
\textcolor{blue}{for opening a terminal in many Linux distributions)}
pyautogui.hotkey('ctrl', 'alt', 't')

\textcolor{purple}{# Action 3}
\textcolor{blue}{# Move the mouse to the terminal and click to focus on the terminal window}
pyautogui.moveTo(260 + 722/2, 103 + 434/2, duration=1)  
pyautogui.click()
\textcolor{blue}{# Type the command to use ffmpeg to trim the video}
pyautogui.write(ffmpeg_command, interval=0.05)
pyautogui.press(`enter')

\textcolor{purple}{# Done, but not follow the demanding in instruction.}
\end{Verbatim}
\end{tcolorbox}

\begin{tcolorbox}[title={Agent > Human Task: \textit{I want to understand the resource usage of my Ubuntu system under normal workloads. Please use the `sar` command in the `sysstat` toolkit to monitor system activity, evaluate the status once every second for 30 seconds, output the results to "System Resources Report.txt" under Desktop.}}]
\begin{Verbatim}[commandchars=\\\{\}]
\textcolor{purple}{# Action 1}
import pyautogui
import time
\textcolor{blue}{# Press the keyboard shortcut to open the terminal}
pyautogui.hotkey(`ctrl', `alt', `t')
time.sleep(1)  
\textcolor{blue}{# Wait a moment for the terminal to open}

\textcolor{purple}{# Action 2}
\textcolor{blue}{# Click inside the terminal to focus on it}
pyautogui.click(600, 300)  
\textcolor{blue}{# Assuming the terminal is in the upper left area of the screen}
\textcolor{blue}{n# Type the command to monitor system activity}
pyautogui.typewrite(`sar 1 30 > ~/Desktop/System Resources Report.txt')

\textcolor{purple}{# Action 3}
\textcolor{blue}{# Click on the `System Resources Report.txt' file to select it}
pyautogui.click(1850, 940)  
\textcolor{blue}{# Coordinates taken from the label position and size}
time.sleep(0.5)
\textcolor{blue}{# Right-click to open the context menu}
pyautogui.click(1850, 940, button=`right')
\textcolor{blue}{# Move to the `Open With Other Application' option in the context menu}
pyautogui.move(0, 70)  
\textcolor{blue}{# Assuming `Open With Other Application' is 70 pixels}
\textcolor{blue}{down from the right-click position}
\textcolor{blue}{# Click on `Open With Other Application'}
pyautogui.click()
time.sleep(0.5)
\textcolor{blue}{# Type `gedit' to select the text editor}
pyautogui.typewrite(`gedit')
\textcolor{blue}{# Press `Enter' to open the file with gedit}
pyautogui.press(`enter')

\textcolor{purple}{# Action 4}
\textcolor{blue}{# Click inside the terminal to ensure it's in focus}
pyautogui.click(500, 200)  
\textcolor{blue}{# Coordinates are estimated to be within the terminal window}
\textcolor{blue}{# Type the command to open the file with gedit}
pyautogui.typewrite(`gedit ~/Desktop/System Resources Report.txt')

\textcolor{purple}{Done!}
\end{Verbatim}
\end{tcolorbox}

\begin{tcolorbox}[title={Agent > Human Task: \textit{Hey, my LibreOffice Writer seems to have frozen and I can't get it to close normally. Can you help me force quit the application from the command line? I'm on Ubuntu and I don't want to restart my computer or lose any other work I have open.}}]

\begin{Verbatim}[commandchars=\\\{\}]
\textcolor{purple}{# Action 1}
import pyautogui
import time
\textcolor{blue}{# Open a terminal window}
pyautogui.hotkey(`ctrl', `alt', `t')
time.sleep(1)
\textcolor{blue}{# Type the command to kill all instances of LibreOffice Writer}
pyautogui.typewrite(`killall soffice.bin')

\textcolor{purple}{# Action 2}
\textcolor{blue}{# Open a terminal window}
pyautogui.hotkey(`ctrl', `alt', `t')
time.sleep(1)
\textcolor{blue}{# Type the command to kill all instances of LibreOffice Writer}
pyautogui.typewrite(`killall soffice.bin')

\textcolor{purple}{# Action 3}
\textcolor{blue}{# Ensure the terminal window is focused}
pyautogui.click(500, 200)  
\textcolor{blue}{# Assuming the terminal is at this position}
\textcolor{blue}{# Type the pkill command to forcefully kill the LibreOffice Writer application}
pyautogui.typewrite(`pkill -9 soffice.bin')

\textcolor{purple}{Done!}
\end{Verbatim}
\end{tcolorbox}


\subsection{Comparative Analysis: Claude-3 vs. GPT-4V}
\label{appendix:claude3}

\begin{figure}[htbp]
    \centering
    \includegraphics[width=\linewidth]{images/qualititive_analysis/claude_error.pdf}
    \caption{
    Screenshots of error example analysis of Claude-3 Opus.
    }
    \label{fig:claude_error}
\end{figure}

\begin{tcolorbox}[title={Claude3 Failed Task1: \textit{I remember there is a file named "secret.docx" on this computer, but I can't remember where it is. Please find the path where this file is stored and copy it to the clipboard.}}]

\begin{Verbatim}[commandchars=\\\{\}]
\textcolor{purple}{# Action 1}
pyautogui.click(tag_10) 
\textcolor{blue}{# click the Files icon to open the file manager}

\textcolor{purple}{# Action 2}
pyautogui.click(tag_21) 
\textcolor{blue}{# click the search icon to open the search bar}

\textcolor{purple}{# Action 3}
pyautogui.typewrite(`secret.docx') 
\textcolor{blue}{# type the file name to search for}
pyautogui.press(`enter') 
\textcolor{blue}{# press enter to start the search}

\textcolor{purple}{# Action 4}
pyautogui.doubleClick(tag_52) 
\textcolor{blue}{# double click the secret.docx file to select it}
pyautogui.hotkey(`ctrl', 'c') 
\textcolor{blue}{# press Ctrl+C to copy the selected file path to clipboard}

\textcolor{purple}{Failed.}
\end{Verbatim}
\end{tcolorbox}


\begin{tcolorbox}[title={Claude3 Failed Task2: \textit{copy the data in Cell B6 in this Libreoffice Calc file and search it in the Chrome browser.}}]

\begin{Verbatim}[commandchars=\\\{\}]
\textcolor{purple}{# Action 1}
import pyautogui
import time
\textcolor{blue}{# Scroll down to view more search results if needed}
pyautogui.press(`pagedown')
time.sleep(1)
pyautogui.press(`pagedown') 
time.sleep(2)
\textcolor{blue}{# The answer appears to be in the snippet under the first search result}
\textcolor{blue}{# `TV 31. u76f4u64adu4e2d.' which translates to `TV 31. Live.'}
\textcolor{blue}{# So channel 31 in Hong Kong is RTHK TV 31}
print(`Found the answer. Channel 31 in Hong Kong is RTHK TV 31.')
\textcolor{blue}{# The original task of copying cell B36 from the spreadsheet}
\textcolor{blue}{and searching it in the browser is complete}
print(`Task complete.')

\textcolor{purple}{# Action 2}
\textcolor{blue}{# Activate the Chrome browser window}
pyautogui.click(tag_3)
time.sleep(1)
\textcolor{blue}{# Click on the address/search bar}
pyautogui.click(tag_30)
\textcolor{blue}{# Paste the copied value `Dewitt'}
pyautogui.hotkey(`ctrl', `v')
\textcolor{blue}{# Press Enter to perform the search}
pyautogui.press(`enter')
print(`Searched for `Dewitt' in the Chrome browser.')

\textcolor{purple}{# Action 3}
print(`DONE')

\textcolor{purple}{# Action 4}
print(`DONE')

\textcolor{purple}{Failed.}
\end{Verbatim}
\end{tcolorbox}

\begin{tcolorbox}[title={Claude3 Failed Task3: \textit{Please help me change all the places in this document that say `text' to `test'.}}]

\begin{Verbatim}[commandchars=\\\{\}]
\textcolor{purple}{# Action 1}
pyautogui.hotkey(`ctrl', `h')
pyautogui.write(`text')
pyautogui.press(`tab')
pyautogui.write(`test')
pyautogui.press(`enter')

\textcolor{purple}{# Action 2}
pyautogui.press(`enter')

\textcolor{purple}{# Action 3}
pyautogui.hotkey(`ctrl', `s')

\textcolor{purple}{Failed.}
\end{Verbatim}
\end{tcolorbox}

\subsection{Analysis on SoM Effectiveness}

\begin{tcolorbox}[title={Task w/o SoM: \textit{turning off the code checking function of the python plugin in vscode}}]

\small
\begin{Verbatim}[commandchars=\\\{\}]
\textcolor{purple}{# Action 1}
import pyautogui
import time
\textcolor{blue}{# Move the mouse to the gear icon (Manage) and click it}
pyautogui.moveTo(472, 1034, duration=0.5) 
pyautogui.click()

\textcolor{purple}{# Action 2}
\textcolor{blue}{# Move the mouse to the \"Settings\" menu item and click it}
\textcolor{blue}{# A rough central point of the `Settings' menu item}
pyautogui.moveTo(497 + 385/2, 832 + 25/2, duration=0.5)  
pyautogui.click()

\textcolor{purple}{# Action 3}
\textcolor{blue}{# Move the mouse to the search settings bar and click it}
\textcolor{blue}{# A rough central point of the search settings bar}
pyautogui.moveTo(783 + 587/2, 390 + 19/2, duration=0.5)  
pyautogui.click()
\textcolor{blue}{# Type the query to search for Python-related settings}
pyautogui.write('python', interval=0.05)

\textcolor{purple}{# Action 4}
\textcolor{blue}{# Move the mouse to the `Edit in settings.json' button and click it}
\textcolor{blue}{# A rough central point of the button}
pyautogui.moveTo(990 + 116/2, 624 + 15/2, duration=0.5)  
pyautogui.click()

\textcolor{purple}{# Action 5}
\textcolor{blue}{# Move the cursor to the correct position}
pyautogui.moveTo(889 + 1/2, 432 + 1/2, duration=0.5)  
pyautogui.click()
\textcolor{blue}{# Type the configuration to disable Python missing imports error reporting}
pyautogui.write(`reportMissingImports: `none'', interval=0.05)
pyautogui.press(`enter')
\textcolor{blue}{# Save the changes}
pyautogui.hotkey(`ctrl', `s')

\textcolor{purple}{# Done}
\end{Verbatim}

\end{tcolorbox}

\begin{tcolorbox}[title={Task w/ SoM: \textit{  turning off the code checking function of the python plugin in vscode}}]
    
\small
\begin{Verbatim}[commandchars=\\\{\}]
\textcolor{purple}{# Action 1}
import pyautogui
import time
pyautogui.hotkey(`ctrl', `shift', `p')

\textcolor{purple}{# Action 2}
pyautogui.write(`Preferences: Open Settings (UI)', interval=0.05)
time.sleep(0.5)
pyautogui.press(`enter')

\textcolor{purple}{# Action 3}
\textcolor{blue}{# Click on the search settings input field}
pyautogui.click(x=600, y=175)  
\textcolor{blue}{# These coordinates may need to be adjusted for your screen resolution}
\textcolor{blue}{# Type `Python' into the search field}
pyautogui.write(`Python', interval=0.05)
pyautogui.press(`enter')

\textcolor{purple}{# Action 4}
\textcolor{blue}{# Click on the search settings bar and type in the query again}
pyautogui.click(x=800, y=150)  
\textcolor{blue}{# Click on the search bar, the coordinates may need adjustment}
pyautogui.typewrite(`Python missing imports', interval=0.25)  

\textcolor{blue}{# ... Can not find the right click area}

\textcolor{purple}{# Failed}
\end{Verbatim}


\end{tcolorbox}


% \subsection{Wrong Calculation of Coordinates}

% \begin{figure}[H]
%     \centering
%     \includegraphics[width=0.8\textwidth]{images/qualititive_analysis/wrong_coordinate_calc/step_1_20240131@151309.png}
%     \caption{multi\_apps/897e3b53-5d4d-444b-85cb-2cdc8a97d903}
%     \label{fig:error_case_multi897e}
% \end{figure}

% \small
% \begin{Verbatim}
% # Wrong action
% import pyautogui
% import time

% # Move the mouse to the "Install Theme" button for the "Dark Fox" theme and click it
% pyautogui.moveTo(945 + 136/2, 2810 + 36/2, duration=1)  # Aim for the center of the button
% time.sleep(0.5)
% pyautogui.click()
% \end{Verbatim}





% \begin{figure}[htbp]
%     \centering
%     \includegraphics[width=\linewidth]{images/qualititive_analysis/wrong_click.pdf}
%     \caption{
%     wrong_click
%     }
%     \label{fig:wrong_click}
% \end{figure}


% \begin{figure}[htbp]
%     \centering
%     \includegraphics[width=\linewidth]{images/qualititive_analysis/cookie_block.pdf}
%     \caption{
%     cookie_block
%     }
%     \label{fig:cookie_block}
% \end{figure}

\subsection{Analysis on SoM Effectiveness}
The SoM, effective for web tasks focusing only on labeled bounding boxes, \textbf{did not} excel in our diverse tasks that required more expertise and flexibility. It shortens the action space of the agent and thus hinders the agent's exploration and adaptability.
For example in Figure \ref{fig:som_analysis}, in the task of "turning off the code checking function of the python plugin." in VS Code, agents without SoM succeeded by editing the settings.json, whereas SoM-equipped agents struggled with finding and deselecting the checkbox in the settings. 
The latter has a longer action path and fails due to incorrect mouse clicks. 
Furthermore, SoM's labeling of elements in professional software also has such errors, that agents without SoM can use a11y tree to calculate and adjust coordinates, while the elements' blocks under SoM can be hard to change during the task.

\begin{figure}[htbp]
    \centering
    \vspace{-5pt}
    \includegraphics[width=\linewidth]{images/qualititive_analysis/som_analysis.pdf}
    \caption{
    A task showcase where the SoM-equipped agent tried to find interactive settings, while the screen-a11ytree-equipped agents directly modified the value in the setting.json file.
    }
    \label{fig:som_analysis}
    \vspace{-10pt}
\end{figure}
